%% =========================================================
%% Chapter: Positive quasi-local constraints from a global gap
%% =========================================================

\chapter{Positive quasi-local constraints from a global gap}%
\label{ch:peps_area_law_constraints}

This chapter follows Section~4 of the two-dimensional area-law manuscript
(Propositions 4.3 and Lemma 4.4). Let $G$ be a finite graph and
$H=\sum_i h_i$ a sum of Hermitian terms, where $h_i$ acts on a support $X_i$
with anchor $a_i\in X_i$, every support has graph diameter at most $R$ and at
most $v_R$ sites, the norms of the terms meeting a site sum to at most $b_0$,
$\|h_i\|\le J$, and the spheres of $G$ grow at most polynomially. Suppose that
$H\Omega=E_0\Omega$ for a unit vector $\Omega$ and that
$H-E_0I\ge\Delta(I-P_\Omega)$ with $\Delta>0$, where
$P_\Omega=|\Omega\rangle\langle\Omega|$. Write $E_{i,l}$ for the normalized
partial-trace expectation onto the graph ball $N_l(a_i)$
(Definition~\ref{def:qc_site_expectation}) and $\alpha=p/(p+1)$ for a fixed
integer $p\ge1$.

The construction filters each term with a smooth spectral filter of width
$\Delta/2$ whose time kernel $f$ decays like $e^{-c|t|^\alpha}$, removes its
ground expectation, and takes the absolute value:
\[
    M_i=\int_{\mathbb R}f(t)e^{itH}h_ie^{-itH}\,dt-\langle\Omega,h_i\Omega\rangle I,
    \qquad k_i=|M_i|/c_*,\qquad c_*=\max\{1,J(\|f\|_1+1)\}.
\]

\begin{lemma}[Order and commutant of the ball expectation]
    \label{lem:al_site_expectation_order}
    \lean{QuantumCircuit.heisenbergEvolution_eq_hermitianUnitaryPath,
        QuantumCircuit.siteExpectation_conjTranspose,
        QuantumCircuit.isHermitian_siteExpectation,
        QuantumCircuit.siteExpectation_nonneg, QuantumCircuit.siteExpectation_mono,
        QuantumCircuit.siteExpectation_le_one,
        QuantumCircuit.siteExpectation_eq_self_of_forall_commute,
        QuantumCircuit.mem_supportedOperators_of_forall_commute,
        Commute.cfcAbs_of_isHermitian,
        QuantumCircuit.siteExpectation_abs_of_mem_supportedOperators}
    \leanok
    \uses{lem:qc_site_expectation_average, lem:qc_site_expectation_properties}
    The expectation $E_K$ onto a set of sites $K$ preserves adjoints and the
    order, maps $[0,I]$ into $[0,I]$, and fixes every operator commuting with
    all operators acting outside $K$; such an operator acts on $K$. In
    particular $E_K$ fixes $|Y|$ for every Hermitian $Y$ acting on $K$.
\end{lemma}
\begin{proof}\leanok
    $E_K$ is the uniform average of the conjugations by products of on-site
    Weyl operators outside $K$, and its values act on $K$.
\end{proof}

\begin{theorem}[Positive replacement of the Hamiltonian]
    \label{thm:al_positive_replacement}
    \lean{TNLean.PEPS.AreaLaw.commute_filterIntegral_of_forall_commute,
        TNLean.PEPS.AreaLaw.positiveConstraint_mem_supportedOperators_component,
        TNLean.PEPS.AreaLaw.localizationPrefactor,
        TNLean.PEPS.AreaLaw.localizationPrefactor_nonneg,
        TNLean.PEPS.AreaLaw.exists_positiveQuasilocalConstraints}
    \leanok
    \uses{thm:qc_graph_ball_localization, lem:al_site_expectation_order,
        lem:qc_lr_closed_region}
    There are constants $C\ge0$ and $c>0$, depending only on
    $p,R,v_R,b_0,J,\Delta$ and the growth of the spheres, such that
    \[
        0\le k_i\le I,\qquad k_i\Omega=0,\qquad
        \sum_ik_i\ge c_*^{-1}(H-E_0I)\ge\frac{\Delta}{c_*}(I-P_\Omega),
    \]
    and for every integer $l\ge0$ the operator $k_{i,l}=E_{i,l}(k_i)$ lies in
    $[0,I]$ and satisfies $\|k_i-k_{i,l}\|\le Ce^{-cl^\alpha}$. Each $k_i$ acts
    on the connected component of $a_i$.
\end{theorem}
\begin{proof}\leanok
    The spectral filter multiplies the matrix element of $h_i$ between energies
    $E$ and $E_0$ by $\chi(E-E_0)$, which vanishes across the gap and equals
    $1$ at $E=E_0$; hence $M_i\Omega=0$, $\sum_iM_i=H-E_0I$ and
    $\|M_i\|\le c_*$. Since $M_i\le|M_i|$ and the kernel of $M_i$ lies in that
    of $|M_i|$, the first display follows. For the tail, the graph-ball
    localization of the dynamics (Theorem~\ref{thm:qc_graph_ball_localization})
    is integrated against $|f|$, splitting the time integral at $|t|$ of order
    $l$; the square-root estimate
    $\||A|-|B|\|\le((\|A\|+\|B\|)\|A-B\|)^{1/2}$ transfers the bound to the
    absolute value, because $E_{i,l}$ fixes $|E_{i,l}M_i|$. Finally the
    Heisenberg evolution of $h_i$ commutes with every operator acting outside
    the component of $a_i$ (Lemma~\ref{lem:qc_lr_closed_region}), hence so do
    the filtered term, its centering and its absolute value.
\end{proof}

\begin{lemma}[Square roots and local channels]
    \label{lem:al_quasilocal_roots}
    \lean{TNLean.PEPS.AreaLaw.sqrt_mul_exp_neg, TNLean.PEPS.AreaLaw.quasilocalRoots}
    \leanok
    \uses{thm:al_positive_replacement, lem:al_site_expectation_order}
    Let $G_i=(I-k_i)^{1/2}$, $K_i=k_i^{1/2}$, $G_{i,l}=(I-k_{i,l})^{1/2}$ and
    $K_{i,l}=k_{i,l}^{1/2}$. These are positive contractions, $G_i\Omega=\Omega$,
    $K_i\Omega=0$, $G_{i,l}^2+K_{i,l}^2=I$, and
    $\|G_i-G_{i,l}\|+\|K_i-K_{i,l}\|\le2\|k_i-k_{i,l}\|^{1/2}
        \le2C^{1/2}e^{-(c/2)l^\alpha}$. The maps
    $\mathcal E_i(B)=G_iBG_i+K_iBK_i$ and
    $\mathcal E_{i,l}(B)=G_{i,l}BG_{i,l}+K_{i,l}BK_{i,l}$ are unital and
    completely positive, contractive after adjoining any auxiliary system, and
    $\|(\mathcal E_i-\mathcal E_{i,l})\otimes\mathrm{id}\|\le4\|k_i-k_{i,l}\|^{1/2}$
    uniformly in the auxiliary dimension. The map $\mathcal E_{i,l}$ fixes every
    operator commuting with all operators acting on $N_l(a_i)$.
\end{lemma}
\begin{proof}\leanok
    The square-root estimate $\|U^{1/2}-V^{1/2}\|\le\|U-V\|^{1/2}$ gives the
    root tails. The map $v\mapsto Gv\oplus Kv$ is an isometry, and
    $\mathcal E(B)$ is the compression of $B\oplus B$ by it, also after
    tensoring with an identity. The difference bound follows from
    $ABA-A'BA'=(A-A')BA+A'B(A-A')$. An operator commuting with the ball algebra
    commutes with $k_{i,l}$, hence with both local roots.
\end{proof}
